\subsection{FINITE SUBSETS OF ORDERED SETS}

PROPOSITION 3. Let E be a right directed preordered set (resp. a lattice, resp. a totally ordered set). Then every non-empty finite subset of E is bounded above (resp. has a least upper bound and a greatest lower bound, resp. has a greatest element and a least element).

The proof is by induction on the number n of elements of the subset under consideration. The result is trivial for n = 1. Let X be a subset of \(n + 1\) elements of E (with \(n \geqslant 1\) ), and put \(X = Y \cup \{x\}\) , where Y has n elements and is therefore not empty. The inductive hypothesis implies the existence of an upper bound (resp. least upper bound, resp. greatest element) y of Y. Since E is right directed (resp. a lattice, resp. totally ordered), \(\{x, y\}\) has an upper bound (resp. least upper bound, resp. greatest element) which is evidently an upper bound (resp. least upper bound, resp. greatest element) of X.

COROLLARY 1. Every totally ordered finite set is well-ordered and has a greatest element.

COROLLARY 2. Every finite ordered set has a maximal element.

For such a set is inductive by Corollary 1 (cf.~§2, no. 4, Theorem 2).

